Ou \todo{Lecture 6} le cote a gauche est 

\begin{align}
 & =\frac{d}{dt}\lim_{ n \to \infty } \sum_{i=1}^{n} \rho(\boldsymbol{r}_{i},t)\Delta V_{i} \\
 & =\lim_{ n \to \infty } \sum_{i=1}^{n} \frac{ \partial \rho }{ \partial t } (\boldsymbol{r}_{i},t)\Delta V_{i} =\iiint _{V}\frac{ \partial \rho }{ \partial t } ~dV 
\end{align}

et a droite

\begin{align}
 & =-\iint _{S}\rho \boldsymbol{v}~d\boldsymbol{S} \\
 & =-\iiint _{V}\boldsymbol{\nabla}(\rho \boldsymbol{v})~dV \\
 & \Rightarrow \iiint _{V}\left( \frac{ \partial \rho }{ \partial t } +\boldsymbol{\nabla}(\rho \boldsymbol{v}) \right)~dV=0
\end{align}
pour n'importe quelle valuer $V$ donc
$$
\frac{ \partial \rho }{ \partial t } +\boldsymbol{\nabla}(\rho \boldsymbol{v})
$$
Ou de maniere equivalent 
$$
\frac{D\rho}{Dt}+\rho(\boldsymbol{\nabla v})=0
$$
\subsection{Equation de continuite (description Lagrangienne)}

Derivation differente avec le meme resultat (voir serie 4.6!!!!).
\begin{remark}
\begin{itemize}
\item Il y a un terme "produit de $\rho$ et $\boldsymbol{v}$" donc l'equation de continuite est non-lineaire.
\item On appelle un fluide incompressible si $\rho$ est constante. Dans ce cas, l'equation de continuite devient:
$$
\boldsymbol{\nabla v}=0
$$
\end{itemize}
\end{remark}

\subsection{Equation d'Euler}
On considere un fluide parfait (pas de frottement interne, donc pas de viscosite, ni de conduction de chaleur). On utilise la description Lagrangienne du fluide (on suit le volume $V$): 

On considere la trajectoire d'elements de fluide d'un instant $t$ a $t+\Delta t$. La quantite de mouvement dans $V(t)$ est
$$
\boldsymbol{P}=\iiint _{V(t)} \rho \boldsymbol{v}~dV\quad
$$
Donc par Newton on a
$$
\frac{d\boldsymbol{P}}{dt}=\sum \boldsymbol{F}^{\mathrm{ext}}
$$
la somme des forces externes agissant sur la partie du fluide dans $V(t)$ qu'on peut calculer

\begin{align}
\frac{d\boldsymbol{P}}{dt} & =\iiint _{V(t)} \frac{D}{Dt}(\rho \boldsymbol{v})+\rho \boldsymbol{v}(\boldsymbol{\nabla v})~dV \\
 & =\underbrace{ \iiint _{V(t)} \rho \boldsymbol{g}~dV }_{ \text{force de gravite} } -\underbrace{ \iint _{S(t)}p~d\boldsymbol{S} }_{ \text{force de pression} }\\
 & \Rightarrow \underbrace{ \frac{D}{Dt}(\rho \boldsymbol{v})+\rho \boldsymbol{v}(\boldsymbol{\nabla v}) }_{ \frac{\rho D\boldsymbol{v}}{Dt}+\underbrace{ \boldsymbol{v}\left( \frac{D\rho}{Dt} +\rho(\boldsymbol{\nabla v})  \right) }_{ =0 }}=\rho \boldsymbol{g}-\boldsymbol{\nabla}p \\
 & \Rightarrow \frac{\rho D\boldsymbol{v}}{Dt}=\rho \boldsymbol{g}-\boldsymbol{\nabla}p
\end{align}

\begin{remark}
Valable strictement si la viscosite est 0. 
\end{remark}
\begin{remark}
L'equation d'Euler n'est pas lineaire donc la resolution en generale est extrêmement complique.
\end{remark}

\subsection{Equation d'etat}

Pour un fluide parfait on a les fonctions $\rho,p,\boldsymbol{v}$ (a cinq composantes) et les equations de continuite et d'Euler. Il nous manque une equation: l'equation d'etat qui depend du type de fluide.

\begin{itemize}
\item Gaz parfait sans echange de chaleur entre elements fluides: 
$$
p\Delta V^{\gamma}=\text{constante le long de la trajectoire}
$$
ou $\gamma=\frac{c_{p}}{c_{v}}$ est l'indice d'adiabicite (le rapport des chaleurs specifiques). Comme $\rho \propto \frac{1}{\Delta V}$ alors 
$$
\frac{D}{Dt}(p\rho^{-\gamma})=0
$$
\item Liquide parfait: Souvent on use l'approximation du fluide incompressible donc $\rho$ est constante et par l'equation de continuite $\boldsymbol{\nabla v}=0$.
\end{itemize}

On a donc un systeme d'equations pour decrire un fluide parfait:
$$
\begin{cases}
\frac{ \partial p }{ \partial t } +\boldsymbol{\nabla}(\rho \boldsymbol{v})=0 \\
\frac{\rho D\boldsymbol{v}}{Dt}-\rho \boldsymbol{g}+\boldsymbol{\nabla}p=0 \\
\text{Equation d'etat}
\end{cases}
$$

Conditions aux limites:
\begin{example}
Si on a des parois immobiles alors $\boldsymbol{v}~d\boldsymbol{S}=0$
\end{example}
